\section{Двойственность в математике. Принцип двойственности. Представление булевых функций в программах.}

\subsection*{Двойственная функция}

\paragraph{Определение}
Пусть $f(x_1, \dots, x_n)$ — булева функция. Функция $f^*$, определяемая как:
$$f^*(x_1, \dots, x_n) \stackrel{\text{def}}{=} \neg f(\neg x_1, \neg x_2, \dots, \neg x_n),$$
называется \textbf{двойственной} к функции $f$.

\paragraph{Свойства}
\begin{enumerate}
    \item \textbf{Инволютивность:} $f^{**} = f$. Это означает, что отношение «быть двойственной» симметрично.
    \item \textbf{Табличное правило:} Чтобы получить таблицу истинности для $f^*$, нужно инвертировать все значения в таблице функции $f$ (и значения аргументов, и значения самой функции) и переставить строки в установленный порядок. Проще говоря: столбец значений $f^*$ читается как инвертированный столбец $f$, записанный снизу вверх.
    \item \textbf{Самодвойственность:} Функция называется самодвойственной, если $f^* = f$. Примеры: $x$, $\neg x$, медиана $M(x, y, z)$. Конъюнкция и дизъюнкция не самодвойственны ($x \land y \leftrightarrow x \lor y$).
\end{enumerate}

\subsection*{Реализация двойственной функции и Принцип двойственности}

\paragraph{Теорема (о реализации)}
Если функция $\varphi$ реализована формулой $\mathcal{J}(f_1, \dots, f_m)$, то двойственная функция $\varphi^*$ реализуется формулой $\mathcal{J}^*(f_1^*, \dots, f_m^*)$, где каждая подфункция заменена на двойственную ей.

\paragraph{Принцип двойственности}
Пусть формула $\mathcal{J}$ в базисе $F = \{f_1, \dots, f_m\}$ реализует функцию $f$. Тогда формула $\mathcal{J}^*$, полученная из $\mathcal{J}$ заменой каждой функции $f_i$ на двойственную ей $f_i^*$, реализует двойственную функцию $f^*$.
\begin{itemize}
    \item \textbf{Следствие:} Если $\mathcal{J}_1 = \mathcal{J}_2$ — равносильность, то $\mathcal{J}_1^* = \mathcal{J}_2^*$ — также равносильность.
    \item \textit{Пример:} Из закона де Моргана $\neg(x \land y) = \neg x \lor \neg y$ по принципу двойственности автоматически следует $\neg(x \lor y) = \neg x \land \neg y$.
\end{itemize}

\subsection*{Представление булевых функций в программах}

В программировании булевы функции представляются несколькими основными способами:

\begin{enumerate}
    \item \textbf{Вектор значений (Truth Table):}
    Функция от $n$ переменных представляется как целое число или массив битов длиной $2^n$. Например, функция $x \land y$ — это вектор $0001_2$ (число 1), а $x \lor y$ — $0111_2$ (число 7). Это самый быстрый способ для функций с малым $n$ (до 5–6 переменных).
    
    \item \textbf{Аналитическое (формульное):}
    Функция хранится как дерево выражений (Expression Tree) или цепочка символов. Используется в системах символьных вычислений и компиляторах.
    
    \item \textbf{Бинарные решающие диаграммы (BDD):}
    Графовое представление, позволяющее компактно хранить и быстро сравнивать функции от большого числа переменных.
    
    \item \textbf{Битовые операции:}
    Процессорные инструкции (AND, OR, NOT, XOR) позволяют вычислять булеву функцию одновременно для 32 или 64 наборов данных, если они упакованы в машинные слова. Это называется \textbf{битовым параллелизмом}.
\end{enumerate}
